\paragraph{未归一化 Wulff 射线的两层平衡谱、仿射压力与中心最终消失}
沿用本小节对均匀凸各向异性体 \(K\) 与 Wulff 形 \(W_K\) 的固定。为考察未归一化尺度轴在有限分辨率上的纯尺度残余，对每个 \(m\ge 1\) 记
\[
F_m:=F_{m+2},
\qquad
N_m:=2^m,
\qquad
N_m=q_mF_m+a_m,
\qquad
0\le a_m<F_m.
\]
再定义未归一化整数尺度 Wulff 射线
\[
\mathfrak W_{K,m}:=\{\rho W_K:\ 0\le \rho\le N_m-1\},
\]
以及 Fibonacci-adic 读出
\[
r_m:\mathfrak W_{K,m}\longrightarrow \ZZ/(F_m\ZZ),
\qquad
r_m(\rho W_K):=[\rho]_{F_m}.
\]
由此诱导的纤维大小函数记为
\[
d_m^{\mathrm{ray}}(r)
:=
\#\{0\le \rho\le N_m-1:\ \rho\equiv r\!\!\pmod{F_m}\}
\qquad
\bigl(r\in \ZZ/(F_m\ZZ)\bigr).
\]
由于该读出只记录整数尺度 \(\rho\) 的模 \(F_m\) 余类，故其有限尺度几何完全退化为区间 \(\{0,\dots,N_m-1\}\) 的余数分配；正因如此，纯尺度 Wulff 射线在纤维统计、压力谱与 gauge 结构上表现出一条刚性的平衡链。

\begin{theorem}[Wulff 射线模 \texorpdfstring{$F_m$}{Fm} 读出的精确两层律与 Fourier 闭式]\label{thm:conclusion-serrin-wulff-ray-two-level-fourier-closed}
\leanverified{paper\_conclusion\_serrin\_wulff\_ray\_two\_level\_fourier\_closed}
对每个 \(m\ge 1\)，恰有 \(a_m\) 个剩余类满足
\[
d_m^{\mathrm{ray}}(r)=q_m+1,
\]
其余 \(F_m-a_m\) 个剩余类满足
\[
d_m^{\mathrm{ray}}(r)=q_m.
\]
在标准代表系 \(r=0,1,\dots,F_m-1\) 下，更精确地有
\[
d_m^{\mathrm{ray}}(r)=q_m+\mathbf 1_{\{0\le r\le a_m-1\}}.
\]
因此对任意非零频率 \(k\in \ZZ/(F_m\ZZ)\setminus\{0\}\)，其离散 Fourier 系数
\[
\widehat d_m^{\mathrm{ray}}(k)
:=
\sum_{r=0}^{F_m-1}d_m^{\mathrm{ray}}(r)e^{2\pi i kr/F_m}
\]
满足显式闭式
\[
\widehat d_m^{\mathrm{ray}}(k)
=
\sum_{r=0}^{a_m-1}e^{2\pi i kr/F_m}
=
e^{\pi i k(a_m-1)/F_m}\,
\frac{\sin(\pi k a_m/F_m)}{\sin(\pi k/F_m)}.
\]
并且有 Parseval 精确预算
\[
\sum_{k\neq 0}\bigl|\widehat d_m^{\mathrm{ray}}(k)\bigr|^2
=
a_m(F_m-a_m).
\]
\end{theorem}
\begin{proof}
由 Euclid 除法
\[
N_m=q_mF_m+a_m
\]
可知，把区间 \(\{0,1,\dots,N_m-1\}\) 按模 \(F_m\) 分到 \(F_m\) 个剩余类时，每个剩余类至少出现 \(q_m\) 次，且恰有 \(a_m\) 个剩余类多出现一次。由于区间从 \(0\) 起始，在标准代表系 \(0,\dots,F_m-1\) 下，额外一次恰落在前 \(a_m\) 个剩余类上，故
\[
d_m^{\mathrm{ray}}(r)=q_m+\mathbf 1_{\{0\le r\le a_m-1\}}.
\]

对 \(k\neq 0\)，有
\[
\sum_{r=0}^{F_m-1}e^{2\pi i kr/F_m}=0,
\]
从而
\[
\widehat d_m^{\mathrm{ray}}(k)
=
q_m\sum_{r=0}^{F_m-1}e^{2\pi i kr/F_m}
+\sum_{r=0}^{a_m-1}e^{2\pi i kr/F_m}
=
\sum_{r=0}^{a_m-1}e^{2\pi i kr/F_m}.
\]
这是一段有限几何级数，故
\[
\widehat d_m^{\mathrm{ray}}(k)
=
\frac{1-e^{2\pi i k a_m/F_m}}{1-e^{2\pi i k/F_m}}
=
e^{\pi i k(a_m-1)/F_m}\,
\frac{\sin(\pi k a_m/F_m)}{\sin(\pi k/F_m)}.
\]

最后，对上述未归一化 Fourier 变换应用 Parseval 恒等式，
\[
\sum_{k=0}^{F_m-1}\bigl|\widehat d_m^{\mathrm{ray}}(k)\bigr|^2
=
F_m\sum_{r=0}^{F_m-1}\bigl(d_m^{\mathrm{ray}}(r)\bigr)^2.
\]
又
\[
\widehat d_m^{\mathrm{ray}}(0)=N_m,
\]
并且
\[
\sum_{r=0}^{F_m-1}\bigl(d_m^{\mathrm{ray}}(r)\bigr)^2
=
(F_m-a_m)q_m^2+a_m(q_m+1)^2
=
F_mq_m^2+2a_mq_m+a_m.
\]
因此
\[
\sum_{k\neq 0}\bigl|\widehat d_m^{\mathrm{ray}}(k)\bigr|^2
=
F_m\bigl(F_mq_m^2+2a_mq_m+a_m\bigr)-(q_mF_m+a_m)^2
=
a_m(F_m-a_m).
\]
证毕。
\end{proof}

\begin{corollary}[纯尺度 Wulff 射线的碰撞率、总变差与 KL 缺口完全闭式]\label{cor:conclusion-serrin-wulff-ray-collision-tv-kl-closed}
定义
\[
p_m^{\mathrm{ray}}(r):=\frac{d_m^{\mathrm{ray}}(r)}{N_m},
\qquad
u_m(r):=\frac1{F_m}.
\]
则碰撞概率满足精确公式
\[
\mathrm{Col}_m^{\mathrm{ray}}
:=
\sum_{r}\bigl(p_m^{\mathrm{ray}}(r)\bigr)^2
=
\frac1{F_m}
+\frac{a_m(F_m-a_m)}{F_mN_m^2}.
\]
特别地，
\[
F_m\,\mathrm{Col}_m^{\mathrm{ray}}
=
1+\frac{a_m(F_m-a_m)}{N_m^2}
\le
1+\frac{F_m^2}{4N_m^2}.
\]
此外，
\[
\bigl\|p_m^{\mathrm{ray}}-u_m\bigr\|_{\mathrm{TV}}
=
\frac{a_m(F_m-a_m)}{F_mN_m},
\]
而 KL 缺口满足闭式
\[
D_{\mathrm{KL}}(p_m^{\mathrm{ray}}\|u_m)
=
\frac{a_m(q_m+1)}{N_m}\log\frac{F_m(q_m+1)}{N_m}
\;+\;
\frac{(F_m-a_m)q_m}{N_m}\log\frac{F_mq_m}{N_m},
\]
并有上界
\[
D_{\mathrm{KL}}(p_m^{\mathrm{ray}}\|u_m)
\le
\log\!\bigl(F_m\,\mathrm{Col}_m^{\mathrm{ray}}\bigr)
\le
\frac{F_m^2}{4N_m^2}
=
O\!\left(\left(\frac{\varphi}{2}\right)^{2m}\right).
\]
\end{corollary}
\begin{proof}
由定理 \ref{thm:conclusion-serrin-wulff-ray-two-level-fourier-closed}，
\[
\sum_r \bigl(d_m^{\mathrm{ray}}(r)\bigr)^2
=
(F_m-a_m)q_m^2+a_m(q_m+1)^2
=
\frac{N_m^2}{F_m}+\frac{a_m(F_m-a_m)}{F_m}.
\]
两边除以 \(N_m^2\) 即得碰撞公式。

对总变差，注意
\[
p_m^{\mathrm{ray}}(r)-u_m(r)=
\begin{cases}
\dfrac{F_m-a_m}{F_mN_m}, & 0\le r\le a_m-1,\\[8pt]
-\dfrac{a_m}{F_mN_m}, & a_m\le r\le F_m-1.
\end{cases}
\]
故
\[
\bigl\|p_m^{\mathrm{ray}}-u_m\bigr\|_{\mathrm{TV}}
=
\frac12\left(
a_m\frac{F_m-a_m}{F_mN_m}
+(F_m-a_m)\frac{a_m}{F_mN_m}
\right)
=
\frac{a_m(F_m-a_m)}{F_mN_m}.
\]

KL 闭式则是两层分布的直接展开。最后，由 Rényi 单调性或 Jensen 不等式，
\[
D_{\mathrm{KL}}(p_m^{\mathrm{ray}}\|u_m)
\le
\log\!\left(F_m\sum_r \bigl(p_m^{\mathrm{ray}}(r)\bigr)^2\right)
=
\log\!\bigl(F_m\,\mathrm{Col}_m^{\mathrm{ray}}\bigr).
\]
再结合前述碰撞公式与
\[
0\le a_m(F_m-a_m)\le \frac{F_m^2}{4},
\]
便得到
\[
\log\!\bigl(F_m\,\mathrm{Col}_m^{\mathrm{ray}}\bigr)
\le
\frac{a_m(F_m-a_m)}{N_m^2}
\le
\frac{F_m^2}{4N_m^2}.
\]
又由于 \(F_m=F_{m+2}\asymp \varphi^m\) 且 \(N_m=2^m\)，故右端是 \(O((\varphi/2)^{2m})\)。证毕。
\end{proof}

\begin{theorem}[未归一化 Wulff 射线的仿射压力律与 escort 均匀化]\label{thm:conclusion-serrin-wulff-ray-affine-pressure-escort-uniformization}
对任意实参数 \(s>0\)，定义纯尺度射线的 \(s\)-阶幂和
\[
S_s^{\mathrm{ray}}(m)
:=
\sum_r \bigl(d_m^{\mathrm{ray}}(r)\bigr)^s
=
a_m(q_m+1)^s+(F_m-a_m)q_m^s.
\]
则极限
\[
P_s^{\mathrm{ray}}
:=
\lim_{m\to\infty}\frac1m\log S_s^{\mathrm{ray}}(m)
\]
存在，并满足严格仿射律
\[
P_s^{\mathrm{ray}}=s\log 2+(1-s)\log\varphi.
\]
进一步，定义 escort 分布
\[
\pi_{m,s}^{\mathrm{ray}}(r)
:=
\frac{\bigl(d_m^{\mathrm{ray}}(r)\bigr)^s}{S_s^{\mathrm{ray}}(m)},
\]
并记
\[
\eta_{m,s}:=\left(1+\frac1{q_m}\right)^s-1.
\]
则有精确公式
\[
\bigl\|\pi_{m,s}^{\mathrm{ray}}-u_m\bigr\|_{\mathrm{TV}}
=
\frac{a_m(F_m-a_m)\eta_{m,s}}
{F_m\bigl(F_m+a_m\eta_{m,s}\bigr)}.
\]
特别地，对每个固定 \(s>0\)，存在常数 \(C_s>0\) 使
\[
\bigl\|\pi_{m,s}^{\mathrm{ray}}-u_m\bigr\|_{\mathrm{TV}}
\le
\frac{C_s}{q_m}
=
O\!\left(\left(\frac{\varphi}{2}\right)^m\right).
\]
因此，纯尺度 Wulff 射线在每个固定 escort 阶上都保持向均匀层收敛；它不产生向顶端纤维层的冻结型集中。
\end{theorem}
\begin{proof}
由 \(F_m=F_{m+2}=\varphi^{m+2}/\sqrt5+O(\varphi^{-m})\)，可得
\[
\frac1m\log q_m\longrightarrow \log\!\left(\frac2\varphi\right).
\]
另一方面，
\[
S_s^{\mathrm{ray}}(m)
=
q_m^s\Bigl(F_m+a_m\eta_{m,s}\Bigr).
\]
由于 \(0\le a_m<F_m\) 且 \(\eta_{m,s}\to 0\)，有
\[
F_m\le F_m+a_m\eta_{m,s}\le F_m(1+\eta_{m,s}),
\]
故
\[
\frac1m\log\bigl(F_m+a_m\eta_{m,s}\bigr)
\longrightarrow
\log\varphi.
\]
于是
\[
\frac1m\log S_s^{\mathrm{ray}}(m)
=
s\,\frac1m\log q_m+\frac1m\log\bigl(F_m+a_m\eta_{m,s}\bigr)
\longrightarrow
s\log\!\left(\frac2\varphi\right)+\log\varphi,
\]
即
\[
P_s^{\mathrm{ray}}=s\log2+(1-s)\log\varphi.
\]

再注意到
\[
(q_m+1)^s=q_m^s(1+\eta_{m,s}),
\qquad
S_s^{\mathrm{ray}}(m)=q_m^s(F_m+a_m\eta_{m,s}),
\]
故 escort 质量在大纤维与小纤维上分别为
\[
\pi_{m,s}^+:=\frac{1+\eta_{m,s}}{F_m+a_m\eta_{m,s}},
\qquad
\pi_{m,s}^-:=\frac{1}{F_m+a_m\eta_{m,s}}.
\]
与均匀分布 \(u_m(r)=1/F_m\) 比较，可得
\[
\pi_{m,s}^+-\frac1{F_m}
=
\frac{\eta_{m,s}(F_m-a_m)}{F_m(F_m+a_m\eta_{m,s})},
\]
\[
\frac1{F_m}-\pi_{m,s}^-
=
\frac{a_m\eta_{m,s}}{F_m(F_m+a_m\eta_{m,s})}.
\]
于是
\[
\bigl\|\pi_{m,s}^{\mathrm{ray}}-u_m\bigr\|_{\mathrm{TV}}
=
\frac12\left(
a_m\frac{\eta_{m,s}(F_m-a_m)}{F_m(F_m+a_m\eta_{m,s})}
+(F_m-a_m)\frac{a_m\eta_{m,s}}{F_m(F_m+a_m\eta_{m,s})}
\right),
\]
即
\[
\bigl\|\pi_{m,s}^{\mathrm{ray}}-u_m\bigr\|_{\mathrm{TV}}
=
\frac{a_m(F_m-a_m)\eta_{m,s}}
{F_m(F_m+a_m\eta_{m,s})}.
\]

最后，对固定 \(s>0\)，由均值定理可知存在常数 \(C_s'>0\) 使
\[
0\le \eta_{m,s}\le \frac{C_s'}{q_m}.
\]
再用
\[
0\le \frac{a_m(F_m-a_m)}{F_m(F_m+a_m\eta_{m,s})}\le \frac14,
\]
便得
\[
\bigl\|\pi_{m,s}^{\mathrm{ray}}-u_m\bigr\|_{\mathrm{TV}}
\le
\frac{C_s}{q_m}
\]
对某个常数 \(C_s>0\) 成立。由于 \(q_m\asymp (2/\varphi)^m\)，故右端为 \(O((\varphi/2)^m)\)。因此 escort 分布在每个固定 \(s>0\) 上都趋于均匀分布，而不是向顶端层集中。证毕。
\end{proof}

\begin{theorem}[纯尺度 Wulff 射线的单比特层宽相变]\label{thm:conclusion-serrin-wulff-ray-onebit-capacity-transition}
定义纯尺度射线在分辨率 \(m\) 上的预言机容量
\[
C_m^{\mathrm{ray}}(B)
:=
\sum_{r\in \ZZ/(F_m\ZZ)}
\min\bigl(d_m^{\mathrm{ray}}(r),2^B\bigr)
\]
及最优成功率
\[
\mathrm{Succ}_m^{\mathrm{ray}}(B):=\frac{C_m^{\mathrm{ray}}(B)}{N_m}.
\]
则对任意整数 \(B\ge 0\)，有精确二相公式
\[
C_m^{\mathrm{ray}}(B)=
\begin{cases}
F_m\,2^B, & 2^B\le q_m,\\[4pt]
N_m, & 2^B\ge q_m+1.
\end{cases}
\]
换言之，从指数小成功率到完全可逆之间只隔着至多一个比特层宽，不存在非平凡中间区间。

更精确地，记
\[
\beta_{\mathrm c}:=\log_2\!\left(\frac2\varphi\right),
\qquad
B_m:=\lfloor \beta m\rfloor.
\]
则有：
\[
\beta<\beta_{\mathrm c}
\quad\Longrightarrow\quad
\mathrm{Succ}_m^{\mathrm{ray}}(B_m)
=
2^{B_m}\frac{F_m}{N_m}
=
\exp\!\bigl(-(\beta_{\mathrm c}-\beta)m\log 2+o(m)\bigr),
\]
\[
\beta>\beta_{\mathrm c}
\quad\Longrightarrow\quad
\mathrm{Succ}_m^{\mathrm{ray}}(B_m)=1
\qquad
\text{当 \(m\) 充分大时}.
\]
\end{theorem}
\begin{proof}
由预言机容量的一般闭式 \ref{thm:fold-oracle-capacity-closed-form}，
\[
C_m^{\mathrm{ray}}(B)
=
\sum_r \min\bigl(d_m^{\mathrm{ray}}(r),2^B\bigr).
\]
再由定理 \ref{thm:conclusion-serrin-wulff-ray-two-level-fourier-closed}，
\[
d_m^{\mathrm{ray}}(r)\in\{q_m,q_m+1\}.
\]
于是
\[
C_m^{\mathrm{ray}}(B)
=
a_m\min(q_m+1,2^B)+(F_m-a_m)\min(q_m,2^B).
\]
由于 \(2^B\) 是整数，而 \(q_m\) 与 \(q_m+1\) 相邻，故只可能出现两种情形：要么 \(2^B\le q_m\)，要么 \(2^B\ge q_m+1\)。前者给出
\[
C_m^{\mathrm{ray}}(B)=F_m\,2^B,
\]
后者给出
\[
C_m^{\mathrm{ray}}(B)=a_m(q_m+1)+(F_m-a_m)q_m=N_m.
\]
这便证明了严格二相律。

再看临界比特率。由 \(q_m\asymp (2/\varphi)^m\)，可写成
\[
q_m=\exp\!\bigl(\beta_{\mathrm c}m\log2+o(m)\bigr).
\]
若 \(\beta<\beta_{\mathrm c}\)，则最终有 \(2^{B_m}\le q_m\)，故
\[
\mathrm{Succ}_m^{\mathrm{ray}}(B_m)
=
\frac{F_m2^{B_m}}{N_m}
=
\exp\!\bigl(-\beta_{\mathrm c}m\log2+\beta m\log2+o(m)\bigr),
\]
即所示指数式。若 \(\beta>\beta_{\mathrm c}\)，则最终有 \(2^{B_m}\ge q_m+1\)，因此
\[
\mathrm{Succ}_m^{\mathrm{ray}}(B_m)=1
\]
当 \(m\) 充分大时成立。证毕。
\end{proof}

\begin{corollary}[黄金 bulk 共振与碰撞爆增不可能由 Wulff 尺度自由度单独产生]\label{cor:conclusion-serrin-wulff-ray-no-golden-bulk-resonance-source}
由推论 \ref{cor:conclusion-serrin-wulff-ray-collision-tv-kl-closed}，
\[
F_m\,\mathrm{Col}_m^{\mathrm{ray}}
=
1+\frac{a_m(F_m-a_m)}{N_m^2}
\longrightarrow 1.
\]
另一方面，定理 \ref{thm:fold-sigma-phi-diverges} 与推论 \ref{cor:fold-collision-scale-diverges} 已给出
\[
\Sigma_\varphi=+\infty,
\qquad
F_{m+2}\mathsf{Col}_m\longrightarrow +\infty.
\]
因此，任何满足
\[
F_{m+2}\mathsf{Col}_m\to+\infty
\]
或更强地承认黄金 bulk 共振闭包发散的观测协议，都不可能由未归一化 Wulff 射线的单参数尺度自由度单独解释；其来源必然包含尺度轴之外的振幅、频域或纤维外置结构。
\end{corollary}
\begin{proof}
由 \(F_m=F_{m+2}\) 与推论 \ref{cor:conclusion-serrin-wulff-ray-collision-tv-kl-closed}，
\[
0\le
F_m\,\mathrm{Col}_m^{\mathrm{ray}}-1
=
\frac{a_m(F_m-a_m)}{N_m^2}
\le
\frac{F_m^2}{4N_m^2}
\longrightarrow 0.
\]
故 \(F_m\,\mathrm{Col}_m^{\mathrm{ray}}\to 1\)。而推论 \ref{cor:fold-collision-scale-diverges} 给出原始 fold 口径下的碰撞尺度无界增长。两种极限互相排斥，故纯尺度 Wulff 射线不可能单独产生黄金 bulk 共振与碰撞爆增。证毕。
\end{proof}

\endinput
